\chapter{THIẾT KẾ GIẢI PHÁP VÀ QUY TRÌNH HỌC MÁY}

\section{Sơ đồ quy trình}
Quy trình gồm các bước: nạp và làm sạch dữ liệu; chia train/test một lần theo phân tầng nhãn; tiền xử lý và tạo đặc trưng
trong Pipeline; tinh chỉnh siêu tham số bằng 5-fold CV trên tập train; đánh giá trên tập test; phân tích lỗi.

\section{Làm sạch dữ liệu}
Các bước không học tham số được thực hiện trước khi chia dữ liệu: bỏ cột định danh, ép kiểu \texttt{TotalCharges} sang số
(ô trống thành thiếu), mã hóa nhãn \texttt{Yes}/\texttt{No} thành 1/0, kiểm tra và loại dòng trùng lặp.

\section{Chia dữ liệu và chống rò rỉ}
Dữ liệu được chia 80\% huấn luyện và 20\% kiểm thử theo phương pháp phân tầng nhãn, với \texttt{random\_state=42}.
Phép chia này được thực hiện \textbf{một lần} và dùng chung cho mọi thí nghiệm, để các mô hình được so sánh trên cùng dữ liệu.

Các bước học tham số được đặt trong \texttt{Pipeline} của scikit-learn: điền trung vị cho biến số, chuẩn hóa \texttt{StandardScaler}
và mã hóa one-hot cho biến phân loại. Các bước này chỉ được fit trên tập huấn luyện, và kiểm thử đơn vị xác nhận điều đó.

\section{Hai luồng tiền xử lý và đặc trưng}
\begin{itemize}
    \item \textbf{Luồng A (cơ bản):} giữ đặc trưng gốc, điền thiếu, chuẩn hóa biến số và one-hot biến phân loại.
    \item \textbf{Luồng B (cải tiến):} bổ sung bốn đặc trưng có lý giải:
    \begin{itemize}
        \item \texttt{num\_services}: số dịch vụ đăng ký; giả thuyết: khách dùng nhiều dịch vụ gắn bó lâu hơn.
        \item \texttt{charge\_per\_tenure}: \texttt{TotalCharges}/(tenure + 1); đo mức chi tiêu trung bình theo thời gian gắn bó.
        \item \texttt{tenure\_group}: nhóm thời gian gắn bó (0--12, 13--24, 25--48, 49+); quan hệ với nguy cơ rời bỏ không tuyến tính.
        \item \texttt{is\_month\_to\_month}: hợp đồng theo tháng, nhóm có nguy cơ cao nhất trong EDA.
    \end{itemize}
\end{itemize}
Các đặc trưng mới chỉ dùng thông tin của chính khách hàng đó, không dùng nhãn nên không gây rò rỉ.

\section{Các thuật toán và siêu tham số}
\begin{itemize}
    \item \textbf{Baseline (mốc so sánh):} luôn dự đoán lớp đa số (không rời bỏ). Mô hình này cho Accuracy cao nhưng không có giá trị dự đoán,
    nên được dùng để kiểm tra xem các mô hình có học được điều gì hay không.
    \item \textbf{K-NN:} số láng giềng $k \in \{5, 11, 21, 41\}$, trọng số đều hoặc theo khoảng cách. Đặc trưng số được chuẩn hóa vì K-NN dựa trên khoảng cách.
    \item \textbf{Naive Bayes Gaussian:} tham số làm mượt phương sai $\in \{10^{-9}, 10^{-7}, 10^{-5}, 10^{-3}\}$.
    \item \textbf{Cây quyết định:} độ sâu tối đa $\in \{3, 5, 7, 10, \text{không giới hạn}\}$ và số mẫu tối thiểu ở lá $\in \{1, 20, 50\}$.
\end{itemize}
Lý do lựa chọn: ba thuật toán nằm trong danh sách đã học của học phần, dễ giải thích và có sự khác biệt về giả định.
Lưới tham số được giữ nhỏ để kết quả có thể tái tạo và giải thích được phạm vi tìm kiếm.

\section{Tiêu chí đánh giá}
Vì nhãn không cân bằng, nhóm báo cáo Accuracy, Precision, Recall, F1-score, ROC-AUC và PR-AUC cùng ma trận nhầm lẫn.
F1 được dùng làm tiêu chí chọn siêu tham số trong CV; PR-AUC đặc biệt phù hợp khi lớp thiểu số là lớp quan tâm.
Mô hình được chọn theo F1 trung bình trong CV trên tập huấn luyện, không chọn theo tập kiểm thử.
